\chapter{ADAMTS13 Activity}

\section*{What Is This Test?}
The ADAMTS13 test measures the activity of a disintegrin and metalloprotease with thrombospondin type 1 motif, member 13 --- the von Willebrand factor (vWF)-cleaving protease. ADAMTS13 is a plasma enzyme that cleaves ultralarge von Willebrand factor (ULvWF) multimers as they are released from endothelial cells and megakaryocytes. When ADAMTS13 activity is severely deficient, ULvWF multimers accumulate, bind platelets, and cause spontaneous platelet aggregation and microvascular thrombosis, producing the clinical syndrome of thrombotic thrombocytopenic purpura (TTP). The test is performed urgently when TTP is suspected --- a clinical diagnosis based on thrombocytopenia and microangiopathic hemolytic anemia (MAHA), often accompanied by neurologic symptoms, renal dysfunction, and fever. An ADAMTS13 activity below 10\% confirms the diagnosis of TTP \cite{joly2017thrombotic}. When the deficiency is caused by autoantibodies inhibiting the enzyme, the disorder is called immune-mediated (acquired) TTP (iTTP); when it is caused by biallelic mutations in the ADAMTS13 gene, it is congenital TTP (cTTP), also called Upshaw-Schulman syndrome. The test is also used to distinguish TTP from other thrombotic microangiopathies (TMA) such as hemolytic uremic syndrome (HUS), and to monitor patients after recovery or during maintenance therapy.

\section*{Alternative Names}
ADAMTS13; von Willebrand Factor-Cleaving Protease Activity; vWF-Cleaving Protease; ADAMTS13 Activity and Inhibitor; ADAMTS13 with Inhibitor Titer; TTP Panel; A Disintegrin and Metalloproteinase with a Thrombospondin Type 1 Motif, Member 13

\section*{Principle}
ADAMTS13 activity is measured by a fluorogenic assay using a recombinant fragment of von Willebrand factor, typically the FRETS-vWF73 peptide, which contains the A2-domain cleavage site of ADAMTS13 (Tyr1605--Met1606). Patient plasma is incubated with the FRET peptide, and the enzyme's cleavage of the peptide releases a fluorescent signal proportional to the ADAMTS13 activity. Results are reported as a percentage of normal pooled plasma (reference range typically 67\%--100\% or 60\%--130\% depending on the laboratory). A mixing study distinguishes immune from congenital TTP: normal plasma is mixed with patient plasma; if ADAMTS13 activity remains low after mixing, an inhibitor (autoantibody) is present. Inhibitor titers are quantified by a Bethesda-like assay (1 Bethesda unit = the amount of antibody that neutralizes 50\% of ADAMTS13 activity in normal plasma). Immunologic assays (ADAMTS13 antigen, ELISA for anti-ADAMTS13 IgG) are available as adjuncts. An activity $<$10\% is diagnostic of TTP, whereas values 10\%--20\% are intermediate and values $>$20\% argue against TTP and favor an alternative TMA. ADAMTS13 is also mildly reduced in liver disease, sepsis, disseminated intravascular coagulation, and pregnancy, so moderate reductions must be interpreted clinically.

\section*{History}
In 1982, Joel Moake and colleagues observed ultralarge von Willebrand factor multimers in the plasma of patients with chronic relapsing TTP, implicating defective vWF processing \cite{moake1982unusually}. In 1996, two independent groups --- Furlan and colleagues in Switzerland and Tsai and colleagues at Einstein College of Medicine --- identified a deficiency of a vWF-cleaving metalloprotease in the plasma of patients with TTP, and distinguished a constitutional deficiency (relapsing, familial) from an autoantibody-mediated deficiency (idiopathic, acquired) \cite{furlan1997deficient, tsai1998antibodies}. The protease was cloned in 2001 and named ADAMTS13 \cite{levy2001mutations}. This discovery transformed TTP from a clinical diagnosis to a laboratory-confirmed one. The FRETS-vWF73 assay, developed by Kokame and colleagues in 2005, became the standard functional test \cite{kokame2005frets}. The clinical trials establishing caplacizumab (a nanobody inhibiting the vWF-platelet interaction) as adjunctive therapy for immune TTP were published in 2016--2019 \cite{scully2019caplacizumab}, and plasma exchange plus corticosteroids remains the cornerstone of acute management, with rituximab for refractory disease.

\section*{Why This Test Is Ordered}
\begin{itemize}
  \item \textbf{Urgent evaluation of suspected thrombotic thrombocytopenic purpura:} In any patient with unexplained thrombocytopenia and microangiopathic hemolytic anemia (schistocytes on smear, elevated LDH, low haptoglobin, elevated indirect bilirubin) --- with or without neurologic symptoms, fever, or renal dysfunction
  \item \textbf{Distinguishing TTP from other thrombotic microangiopathies:} Hemolytic uremic syndrome (typically diarrhea-associated, Shiga toxin-mediated), atypical HUS (complement-mediated), drug-induced TMA, and DIC
  \item \textbf{Confirmation of immune-mediated (acquired) TTP:} With inhibitor testing to demonstrate ADAMTS13 autoantibodies
  \item \textbf{Diagnosis of congenital TTP (Upshaw-Schulman syndrome):} In neonates, children, and adults with recurrent TTP, in pregnancy (especially third trimester), or with a family history
  \item \textbf{Monitoring after TTP:} Serial ADAMTS13 activity and inhibitor titers to guide maintenance therapy (rituximab, immunosuppression) and predict relapse
  \item \textbf{Evaluation of unexplained thrombocytopenia in pregnancy:} Where TTP must be distinguished from HELLP syndrome and preeclampsia
\end{itemize}

\section*{Primary vs Secondary Test}
ADAMTS13 activity is the primary diagnostic test for TTP but is ordered as an urgent reflex when the clinical constellation of thrombocytopenia and MAHA is present. Because results may take 1--3 days (send-out testing), plasma exchange is started empirically before the result returns whenever TTP is clinically suspected --- treatment must never be delayed awaiting the ADAMTS13 result.

\section*{How This Test Is Performed}
A venous blood sample is collected in a blue-top (sodium citrate) tube, typically 3--5 mL. Because ADAMTS13 is labile, the sample should be transported on ice and processed promptly (centrifuged and frozen at --80\textdegree{}C if the assay will be delayed). The test is almost always performed at a reference laboratory using the FRETS-vWF73 fluorescence assay. The activity is reported as a percentage, and when activity is low, an inhibitor mixing study and titer are added. Results are typically available within 24--72 hours, depending on laboratory transport. The CBC, peripheral blood smear (for schistocytes), LDH, haptoglobin, bilirubin, and creatinine are drawn at the same time to document the TMA.

\section*{What Preparation Is Needed}
No fasting is required. The patient should inform the healthcare provider of: any recent plasma transfusion or plasma exchange (which can transiently raise ADAMTS13 activity and mask a low result --- ideally the sample is drawn before plasma exchange), recent medications (quinine, ticlopidine, and certain chemotherapies can cause drug-induced TMA), pregnancy, liver disease, sepsis, and any personal or family history of TTP or unexplained anemia/thrombocytopenia.

\section*{Contraindications and Precautions}
\begin{itemize}
  \item There are no absolute contraindications to the blood draw. Important interpretive cautions:
  \item (1) an ADAMTS13 activity $<$10\% confirms TTP, but a normal or moderately reduced activity ($>$20\%) does not exclude all TMAs --- it redirects the diagnosis to HUS, atypical HUS, or other causes;
  \item (2) ADAMTS13 is modestly reduced in liver disease, sepsis, DIC, and pregnancy, so intermediate values (10\%--20\%) must be interpreted in context;
  \item (3) a sample drawn after plasma exchange or plasma infusion may show falsely normalized activity;
  \item (4) platelet transfusion is generally contraindicated in TTP (except for life-threatening bleeding) because it may fuel thrombosis --- this precaution applies regardless of the test result;
  \item (5) congenital TTP may present for the first time in adulthood (especially in pregnancy) despite a lifelong mutation;
  \item (6) the inhibitor mixing study can be falsely negative in low-titer cases, and the ADAMTS13 antigen assay may help;
  \item (7) the diagnosis is clinical and must prompt plasma exchange immediately, before results return.
\end{itemize}
\section*{Risks}
Risks are those of routine venipuncture: mild pain, bruising, small hematoma at the collection site, lightheadedness or vasovagal syncope, and extremely rarely, infection or nerve injury. The test itself carries no additional risk; the urgent treatment (plasma exchange) that follows is a separate procedure with its own risks, including catheter complications, transfusion reactions, and citrate toxicity.

\section*{What Happens After the Test}
If ADAMTS13 activity is $<$10\%, the diagnosis of TTP is confirmed. In immune TTP (inhibitor present, most adult cases), the patient is treated with daily therapeutic plasma exchange, high-dose corticosteroids, and caplacizumab (a vWF-binding nanobody), with rituximab for refractory disease. In congenital TTP (no inhibitor, biallelic mutation), plasma infusion alone is sufficient to replace the missing enzyme. Patients who survive the acute episode require long-term follow-up with serial ADAMTS13 activity and inhibitor titers, since relapse risk is elevated when activity falls or inhibitor titers rise. If ADAMTS13 is $>$20\%, an alternative TMA is pursued (Shiga toxin testing for STEC-HUS, complement gene panel for atypical HUS, drug review, ADAMTS13 in pregnancy). All patients are followed for chronic kidney disease, neurocognitive deficits, and, in congenital TTP, ongoing replacement therapy.

\section*{Understanding Results}

\subsection*{Below Normal (Severe Deficiency, $<$10\%)}
\begin{itemize}
  \item \textbf{Immune-mediated (acquired) TTP:} Autoantibodies inhibit ADAMTS13; accounts for more than 90\% of adult TTP; associated with autoimmune disease, pregnancy, and certain drugs; confirmed by a positive inhibitor
  \item \textbf{Congenital TTP (Upshaw-Schulman syndrome):} Biallelic ADAMTS13 mutations cause lifelong deficiency; presents in neonates, children, or adults (often in pregnancy or with infection/stress); no inhibitor is present
  \item \textbf{Pregnancy-associated TTP:} Congenital TTP frequently first manifests in the third trimester; distinguished from HELLP syndrome by ADAMTS13 activity
\end{itemize}

\subsection*{Moderately Reduced (10\%--67\%)}
\begin{itemize}
  \item \textbf{Intermediate (10\%--20\%):} Overlaps with TTP; clinical judgment required; often treated as TTP until results clarify
  \item \textbf{Mildly reduced:} Liver disease, sepsis, disseminated intravascular coagulation, pregnancy, or recent surgery; not diagnostic of TTP
\end{itemize}

\subsection*{Normal ($>$67\% or above the laboratory cutoff)}
\begin{itemize}
  \item \textbf{Alternative TMA:} Hemolytic uremic syndrome (STEC-HUS), atypical HUS (complement-mediated), drug-induced TMA, HELLP syndrome, malignant hypertension, or autoimmune hemolysis
  \item \textbf{After plasma exchange:} Transiently normalized activity from infused normal plasma
\end{itemize}

\section*{Normal Results}
\begin{itemize}
  \item \textbf{ADAMTS13 activity:} 67\%--100\% of normal pooled plasma in most healthy adults (some laboratories report 60\%--130\%)
  \item \textbf{Diagnostic threshold for TTP:} Activity $<$10\% is confirmatory; 10\%--20\% is intermediate; $>$20\% is generally inconsistent with TTP
  \item \textbf{Inhibitor titer:} Absent (negative) in normal plasma and in congenital TTP; present in immune-mediated TTP
  \item \textbf{ADAMTS13 antigen:} Normal or low; not diagnostic by itself
\end{itemize}

\section*{Follow-up Tests}
\begin{itemize}
  \item \textbf{Complete blood count with platelet count and peripheral blood smear:} For schistocytes and to follow platelet recovery
  \item \textbf{Lactate dehydrogenase, haptoglobin, indirect bilirubin, reticulocyte count:} To document and follow microangiopathic hemolysis
  \item \textbf{Serum creatinine and urinalysis:} To assess renal involvement
  \item \textbf{Inhibitor titer and repeat ADAMTS13 activity:} To confirm immune TTP and guide maintenance therapy
  \item \textbf{Shiga toxin stool testing:} To confirm or exclude STEC-HUS
  \item \textbf{Complement gene panel (atypical HUS):} When ADAMTS13 is normal and HUS is suspected
  \item \textbf{ADAMTS13 gene sequencing:} To confirm congenital TTP
  \item \textbf{Direct antiglobulin test (Coombs test):} To exclude autoimmune hemolytic anemia
  \item \textbf{ADAMTS13 antigen and anti-ADAMTS13 antibody assays:} Adjunctive testing in ambiguous cases
\end{itemize}

\section*{References}
\bibliographystyle{unsrtnat}
\bibliography{references}

\clearpage
\section*{Regulatory Status and Authorization Date}
This chapter describes a test category, not a specific commercial product. FDA, CDSCO, EU IVDR, and MHRA approval or registration dates therefore cannot be assigned without the assay manufacturer, product name, instrument, and intended use. For laboratory-developed tests, an individual agency approval date may not exist. See the Preface for the interpretation of regulatory status.

\clearpage
